\medskip

\begin{enumerate}
\item  On a, pour tout $x\in[0~;~1]$ : $f(x)=6x^3+6x^2$.

Donc :
\begin{align*}
\int_0^1 f(x)\mathrm{d}x &= \int_0^1\left(6x^3+6x^2\right)\mathrm{d}x \\
&= \left[\frac32x^4+2x^3\right]_0^1 \\
&= \left[\frac32\times1^4+2\times1^3\right] - \left[\frac32\times0^4+2\times0^3\right] \\
&=\frac32+2 \\
&= \frac72.
\end{align*}

\item Comme la fonction $f$ est positive sur $[0~;~1]$, l'intégrale calculée représente l'aire, en unité d'aire, du domaine situé sous la courbe $\mathcal{C}_2$ entre les droites d'équations $x=0$ et $x=1$.

Dans le contexte de l'exercice, cette aire correspond au domaine grisé $\mathcal{D}_1$ et est donc égale à 3,5 unités d'aire.

\item  Le rectangle OABC a pour largeur $1$ et pour hauteur $12$, donc son aire vaut $1\times 12=12$ unités d'aire.

Les trois domaines $\mathcal{D}_1$, $\mathcal{D}_2$ et $\mathcal{D}_3$ forment tout le rectangle OABC. Ainsi :
\begin{align*}
\mathcal{A}(\mathcal{D}_3) &= 12-\mathcal{A}(\mathcal{D}_1)-\mathcal{A}(\mathcal{D}_2) \\
&= 12-\dfrac72-\dfrac{135}{32} \\
&= \dfrac{384}{32}-\dfrac{112}{32}-\dfrac{135}{32} \\
&= \dfrac{137}{32} \\
&\approx \np{4,28} \text{ unités d'aire.}
\end{align*}
\end{enumerate}

\bigskip
